\origpage{44}{}
\textbf{Proof.} By the main result of [14], there is a linked tree-decomposition $(T, W)$ of $G$, with width $\le w$\fix{, such that no edge of $G$ belongs to $W(t)$ for two distinct $t \in V(T)$ (from any linked tree-decomposition, delete each edge of $G$ from all but one of the graphs $W(t)$ that contain it; this keeps the sets $W(e)$)}.\fixnote{fix:11.2-proof}{Without this property the graphs $W'(t_1), W'(t_2)$ below need not give a tree-decomposition. If an edge $f$ of $W(t_0)$ is also in graphs $W(t)$ on both sides of $e_0$, then it is in only one of $W'(t_1), W'(t_2)$, and the path of $T'$ between two vertices whose graphs contain $f$ passes through the other one. For example, let $G$ be one edge $f$, let $T$ be a star with center $t_0$ and four leaves, and let every $W(t)$ be $G$. With the property, $f$ is in no $W(t)$ with $t \ne t_0$, and $(T', W')$ has the property too.} Choose it such that $\Sigma(d(t) - 3)$ is minimum, where $d(t)$ is the degree of $t$ in $T$ and the sum is over all $t \in V(T)$ with degree $\ge 3$. Suppose that some $t_0 \in V(T)$ has degree $\ge 4$. Then we may assume that there is a tree $T'$ with an edge $e_0$ with ends $t_1, t_2$, such that $T$ is obtained from $T'$ by contracting $e_0$, and $t_0$ results from identifying $t_1$ and $t_2$, and such that $t_1$ and $t_2$ both have degree $\ge 3$. Let $W'(t) = W(t)$ if $t \in V(T') \setminus \{t_1, t_2\}$, and let $W'(t_1), W'(t_2)$ be graphs both with vertex set $W(t_0)$, so that $W'(t_1) \cup W'(t_2) = W(t_0)$ and $E(W'(t_1) \cap W'(t_2)) = \emptyset$. Then $(T', W')$ is a tree-decomposition of $G$ with width $\le w$, and $W'(e_0) = V(W(t_0))$, and if $e \in E(T') \setminus \{e_0\}$ then $e \in E(T)$ and $W'(e) = W(e)$. Let $e_1, e_2 \in E(T')$ with $|W'(e_1)| = |W'(e_2)| = k$ say, and $|W'(e)| \ge k$ for all edges $e$ of the path of $T'$ between $e_1$ and $e_2$. We claim that there are $k$ mutually vertex-disjoint paths between $W'(e_1)$ and $W'(e_2)$. If $e_1, e_2 \ne e_0$, then $e_1, e_2 \in E(T)$ and the claim therefore holds, since $W'(e_{\delm{2}\fix{i}}) = W(e_i)$ ($i = 1, 2$)\fixnote{fix:11.2-proof-linked}{The reason is $W'(e_i) = W(e_i)$ for $i = 1, 2$, which the proof showed a line earlier for every edge other than $e_0$.} and $(T, W)$ is linked. If $e_1 = e_2 = e_0$ the claim is trivial. If $e_1 = e_0$ and $e_2 \ne e_0$, let $e_3$ be the edge of $T$ different from $e_1$, between $e_1$ and $e_2$ with a common end $t_1$ say with $e_1$. Then $|W(e_3)| \ge k$, but
\begin{equation*}
W(e_3) \subseteq V(W(t_1)) = W(e_0)
\end{equation*}
and $|W(e_0)| = k$; and so $W(e_0) = W(e_3)$. Since there are $k$ mutually disjoint paths of $G$ between $W(e_3) = W(e_0)$ and $W(e_2)$, the claim follows. Consequently $(T', W')$ is linked, contrary to the choice of $(T, W)$. We deduce that there is no such $t_0$, and so \del{(i) holds}\fix{every vertex of $T$ has degree $\le 3$}, as required.\fixnote{fix:11.2-proof-end}{The statement of (11.2) has no (i). What is meant is its conclusion, that every vertex of $T$ has degree $\le 3$.} $\square$

Secondly, we need the following.

\textbf{11.3} \textit{Let $(T, W)$ be a tree-decomposition of $G$, and let $e_1, e_2 \in E(T)$ be distinct. Let $T_0, T_1, T_2$ be the three components of $T \setminus e_1, e_2$, where $e_i$ has an end in $V(T_0)$ and in $V(T_i)$ ($i = 1, 2$). Suppose that $W(T_0) \subseteq W(T_1) \cap W(T_2)$. Let $T'$ be the tree obtained from $T_1 \cup T_2$ by adding an edge joining the end of $e_1$ in $T_1$ to the end of $e_2$ in $T_2$; and let $W'$ be the restriction of $W$ to $V(T')$. Then $(T', W')$ is a tree-decomposition of $G$, and if $(T, W)$ is linked then so is $(T', W')$.}

\textbf{Proof.} It is straightforward to check that $(T', W')$ is a tree-decomposition of $G$. \fix{To see this, let $e_i$ have ends $a_i \in V(T_i)$ and $c_i \in V(T_0)$ ($i = 1, 2$). Since $T_1, T_2$ are disjoint trees and the new edge joins $a_1$ to $a_2$, $T'$ is a tree. Let $x, z \in V(T')$. If $x, z$ are both in $T_1$ or both in $T_2$, the path of $T$ between them lies in that tree and is also their path in $T'$. Otherwise the path of $T$ between them passes through $a_1, c_1, c_2, a_2$ (since $e_1, e_2$ are the only edges between $T_0$ and $T_1 \cup T_2$), and their path in $T'$ is obtained from it by replacing its subpath between $a_1$ and $a_2$ by the new edge. In both cases every vertex $y$ of the path of $T'$ between $x$ and $z$ is on the path of $T$ between them, and so $W'(x) \cap W'(z) \subseteq W'(y)$. Since $W(T_0) \subseteq W(T_1) \cap W(T_2)$, every vertex and edge of $G$ belongs to $W(T_1) \cup W(T_2)$. Hence $(T', W')$ is a tree-decomposition of $G$.} Let $X = V(W(T_1) \cap W(T_2))$. Since $W(T_0) \subseteq W(T_1) \cap W(T_2)$ it follows that $W(e) = X$ for every edge $e$ of the path of $T$ between $e_1$ and $e_2$, and now the second claim follows easily. \fix{To see this, let $f \in E(T')$. If $f \ne a_1 a_2$ then $f \in E(T)$, and $W'(f) = W(f)$, since both are the vertex set of the intersection of the graphs at the two ends of $f$; let $\bar{f} = f$. Since $T_1, T_2$ are the components of $T' \setminus a_1 a_2$, we have $W'(a_1 a_2) = X = W(e_1)$; let $\bar{f} = e_1$ if $f = a_1 a_2$. Now let $f_1, f_2 \in E(T')$ be distinct (if $f_1 = f_2$ there is nothing to prove), with $|W'(f_1)| = |W'(f_2)| = k$, such that $|W'(f)| \ge k$ for all edges $f$ of the path $R'$ of $T'$ with first edge $f_1$ and last edge $f_2$. Let $R$ be the path of $T$ with first edge $\bar{f}_1$ and last edge $\bar{f}_2$, and let $e \in E(R)$. If $e \in E(T_1 \cup T_2)$, then $e \in E(R')$ by the description of paths above, and so $|W(e)| = |W'(e)| \ge k$. Otherwise $e$ is on the path of $T$ between $a_1$ and $a_2$, so $W(e) = X$, and $a_1 a_2 \in E(R')$ (since either $a_1 a_2$ is $f_1$ or $f_2$, or $R$ passes from $T_1$ to $T_2$); hence $|W(e)| = |W'(a_1 a_2)| \ge k$. Since $(T, W)$ is linked, there are $k$ mutually vertex-disjoint paths of $G$ between $W(\bar{f}_1) = W'(f_1)$ and $W(\bar{f}_2) = W'(f_2)$. Hence $(T', W')$ is linked.}\fixnote{fix:11.3-proof-added}{The paper leaves both checks to the reader. The arguments above are the ones used in the formalization.} $\square$

\textbf{Proof of (11.1).} Choose $w \ge 0$ so that (1.1) is satisfied. Choose $n \ge 0$ so that (2.7) is satisfied with $k$ replaced by $w + 1$, and let $r(p) = 2^{n-1}w(w + 1)$. We claim that (11.1) holds. For we prove
